\chapter{LLM Training Systems}

\begin{itemize}
    \item \term{Goal}: understand the exact execution loop, memory bill, precision tricks, and multi-GPU partitioning behind training a modern LLM.
    \item \term{Running examples}:
    \begin{itemize}
        \item \term{Qwen2.5-72B}: clean dense + GQA baseline.
        \item \term{DeepSeek-V3}: 671B total / 37B active MoE + MLA + FP8 + DualPipe.
    \end{itemize}
    \item \term{Runnable companion}: \code{01-training/code/tiny\_train.py} --- standalone $\sim$280-line PyTorch trainer using the exact same block structure, runnable on NVIDIA H100 (default) or CPU.
\end{itemize}

% -----------------------------------------------------------------------------
\section{The Core Loop: One Objective, Four Phases}

\begin{itemize}
    \item \term{Input}: text converted by a tokenizer into integer IDs of shape \shape{[B, T]}.
    \begin{itemize}
        \item \code{B} = batch size (sequences processed together).
        \item \code{T} = sequence length (tokens per sequence, typically \code{4,096}--\code{131,072}).
        \item Vocabulary size \code{V}: \code{152,064} (Qwen2.5), \code{129,280} (DeepSeek-V3), \code{256} (our byte-level toy model).
    \end{itemize}
    \item \term{Self-supervised task (next-token prediction)}:
    \begin{itemize}
        \item Model input \code{x = tokens[:, :-1]} (positions \code{0..T-2}).
        \item Ground-truth target \code{y = tokens[:, 1:]} (positions \code{1..T-1}).
        \item One sequence of length \code{T} produces \code{T - 1} supervised predictions simultaneously.
    \end{itemize}
    \item \term{One training step}:
    \begin{enumerate}
        \item \term{Forward}: map \shape{[B, T]} token IDs $\rightarrow$ \shape{[B, T, V]} unnormalized scores (logits); retain intermediate activations needed for backprop.
        \item \term{Loss}: compare predicted distribution at every \code{(b, t)} against true next token \code{y[b, t]} (cross-entropy $\rightarrow$ 1 scalar).
        \item \term{Backward}: walk the graph in reverse to compute one gradient tensor per weight tensor.
        \item \term{Optimizer step (AdamW)}: update every weight using its gradient and running statistics (\code{m}, \code{v}), then zero gradients.
    \end{enumerate}
    \item \term{Core systems property of training}:
    \begin{itemize}
        \item All \code{T} input tokens are known upfront.
        \item Every linear projection multiplies a \shape{[B*T, D\_in]} matrix by a \shape{[D\_in, D\_out]} weight matrix $\rightarrow$ large dense GEMMs.
        \item Result: training is overwhelmingly \term{compute-bound} (Tensor Cores stay busy unless memory or collectives stall them).
        \item \emph{Contrast with Chapter 2}: autoregressive decode generates 1 token at a time (\code{T = 1}) $\rightarrow$ tiny GEMVs $\rightarrow$ \term{memory-bandwidth-bound}.
    \end{itemize}
\end{itemize}

% -----------------------------------------------------------------------------
\section{Inside One Transformer Block (Tensor Shapes Only)}

\noindent\begin{minipage}{\linewidth}
Modern LLMs (Qwen, DeepSeek, Llama) share the same pre-norm decoder-only skeleton:
\begin{lstlisting}
tokens [B, T]
  --> Token Embedding ---------------------------------> x [B, T, D]
        --> Repeat L layers:
              |-- RMSNorm --> Attention (token mixing) ----> + residual
              |-- RMSNorm --> FFN / MoE (channel mixing) --> + residual
        --> Final RMSNorm --> LM Head -----------------> logits [B, T, V]
\end{lstlisting}
\end{minipage}

\begin{itemize}
    \item \term{RMSNorm}: rescales each token's \code{D}-element vector to unit root-mean-square; shape stays \shape{[B, T, D]}. Cheaper than LayerNorm (no mean subtraction, fewer memory reads).
    \item \term{Attention (how tokens talk to earlier tokens)}:
    \begin{itemize}
        \item Linear projections form \code{Q} \shape{[B, H\_q, T, d\_h]}, \code{K} \shape{[B, H\_kv, T, d\_h]}, \code{V} \shape{[B, H\_kv, T, d\_h]}.
        \item \term{RoPE}: injects position by rotating adjacent channel pairs inside \code{Q} and \code{K} based on index \code{t}; zero extra embedding parameters.
        \item \term{Causal mask}: token \code{t} may only attend to positions \code{0..t} (blocks peeking at the target \code{t+1}).
        \item Output projection merges all heads back to \shape{[B, T, D]}.
        \item \term{Three attention flavors (directly sets Chapter 2's KV cache size)}:
        \begin{itemize}
            \item \term{MHA}: \code{H\_kv = H\_q} (largest KV state).
            \item \term{GQA (Qwen2.5)}: \code{H\_kv < H\_q} (multiple Q heads share one KV head).
            \item \term{MLA (DeepSeek-V3)}: compresses K and V jointly into a small latent vector (\code{d\_c = 512}) + decoupled RoPE key (\code{d\_r = 64}), decompresses on the fly during attention.
        \end{itemize}
    \end{itemize}
    \item \term{FFN / MoE (where most parameters live; tokens processed independently)}:
    \begin{itemize}
        \item \term{Dense SwiGLU (Qwen2.5)}: two up-projections (\code{gate}, \code{up}) expand \shape{[B, T, D]} $\rightarrow$ \shape{[B, T, D\_ff]}; elementwise \code{silu(gate) * up}; down-project back to \shape{[B, T, D]}. Every token touches 100\% of FFN weights.
        \item \term{Sparse DeepSeekMoE (DeepSeek-V3, layers 4--61)}: router projects each token \shape{[D]} $\rightarrow$ scores over \code{256} routed experts and selects top \code{8}, plus \code{1} always-on shared expert. Each expert is a small SwiGLU (\code{D\_ff = 2048}) $\rightarrow$ only \code{37B / 671B} (\term{5.5\%}) of weights execute per token.
    \end{itemize}
\end{itemize}

\par\vspace{4pt}\noindent\begin{minipage}{\linewidth}
\small
\begin{tabularx}{\linewidth}{@{}l Y Y Y@{}}
\toprule
\term{Spec} & \code{tiny\_train.py} & \term{Qwen2.5-72B (Dense)} & \term{DeepSeek-V3 (MoE)} \\
\midrule
Total / Active params & 459K / 459K & 72.7B / 72.7B & 671B / 37B \\
Layers (\code{L}) & 2 & 80 & 61 (3 dense + 58 MoE) \\
Hidden dim (\code{D}) & 128 & 8,192 & 7,168 \\
Attention & GQA (4 Q / 2 KV) & GQA (64 Q / 8 KV) & MLA (128 heads, latent 512+64) \\
FFN / Experts & SwiGLU (\code{D\_ff=384}) & SwiGLU (\code{D\_ff=29,568}) & 1 shared + 256 routed (top-8, \code{D\_ff=2,048}) \\
Vocab (\code{V}) & 256 (UTF-8 bytes) & 152,064 & 129,280 \\
\bottomrule
\end{tabularx}
\end{minipage}\par\vspace{4pt}

% -----------------------------------------------------------------------------
\section{Toy Example Walkthrough (\code{code/tiny\_train.py})}

\begin{itemize}
    \item \term{What it builds}:
    \begin{itemize}
        \item Raw UTF-8 byte tokenizer (\code{V = 256}), zero external tokenizer files.
        \item 2-layer pre-norm Transformer with \code{RMSNorm}, \code{RoPE}, \code{GQA} causal attention (via PyTorch SDPA, which dispatches FlashAttention on H100), and \code{SwiGLU}.
        \item 250 AdamW steps on an embedded systems corpus ($\sim$3s on H100, $\sim$16s on CPU).
    \end{itemize}
    \item \term{What to verify in stdout}:
    \begin{enumerate}
        \item \term{Step-0 shape trace}: \shape{(32, 64)} $\rightarrow$ \shape{(32, 64, 128)} $\rightarrow$ \code{Q} \shape{(32, 4, 64, 32)} / \code{K,V} \shape{(32, 2, 64, 32)} $\rightarrow$ logits \shape{(32, 64, 256)}.
        \item \term{Memory accounting}: theoretical static state is \code{7.35 MB} (AdamW takes \code{5.51 MB} alone, 6$\times$ the BF16 weights), while peak GPU HBM reaches \code{108.72 MB} once activations and workspaces are allocated.
        \item \term{Convergence}: loss starts near \code{5.55} ($\ln 256$, uniform random guess) and drops to \code{0.06}, producing an exact greedy continuation.
    \end{enumerate}
\end{itemize}

\begin{lstlisting}
# Default: CUDA GPU (tested on NVIDIA H100) with BF16 autocast
python 01-training/code/tiny_train.py

# Force CPU execution (FP32)
python 01-training/code/tiny_train.py --device cpu
\end{lstlisting}

% -----------------------------------------------------------------------------
\section{The Memory Bill: Why a Single H100 Runs Out of HBM Fast}

Training HBM divides into \term{static state} (fixed by parameter count \code{P}) and \term{dynamic activations} (scales with batch size and context length).

\subsection{Static State: The 16-Byte Rule (Mixed Precision + AdamW)}

\par\vspace{4pt}\noindent\begin{minipage}{\linewidth}
\small
\begin{tabularx}{\linewidth}{@{}l l c Y@{}}
\toprule
\term{Component} & \term{Precision} & \term{Bytes / Param} & \term{Role} \\
\midrule
Model weights & BF16 & 2 & Forward and backward GEMMs \\
Gradients & BF16 & 2 & Output of backward pass \\
Master weights & FP32 & 4 & Small updates underflow in BF16 \\
AdamW 1st moment (\code{m}) & FP32 & 4 & Running mean of gradients \\
AdamW 2nd moment (\code{v}) & FP32 & 4 & Running squared gradients \\
\midrule
\term{Total static floor} & --- & \term{16 $\times$ P} & \term{Before batch size = 1} \\
\bottomrule
\end{tabularx}
\end{minipage}\par\vspace{4pt}

\begin{itemize}
    \item \term{Reality check on 1$\times$ NVIDIA H100 (94 GB HBM)}:
    \begin{itemize}
        \item \term{7B dense model}: \code{16 $\times$ 7B = 112 GB} $\rightarrow$ already OOMs a single 94 GB H100 with zero activations.
        \item \term{Qwen2.5-72B}: \code{16 $\times$ 72.7B $\approx$ 1.16 TB} static state ($\ge$13--16 H100s minimum).
        \item \term{DeepSeek-V3 (671B)}: \code{16 $\times$ 671B $\approx$ 10.7 TB} static state ($>$120 H100s minimum).
    \end{itemize}
    \item \term{Takeaway}: \emph{optimizer states (12 bytes/param), not model weights (2 bytes/param), dominate training HBM.}
\end{itemize}

\subsection{Dynamic Activations (The Context-Length Spike)}

\begin{itemize}
    \item Backward pass requires intermediate forward tensors (inputs to every linear layer, norm, and activation).
    \item Scales roughly with \code{B $\times$ T $\times$ L $\times$ D} --- explodes rapidly at long context (\code{T = 32K--128K}).
    \item \term{Two mandatory fixes in every real LLM trainer}:
    \begin{enumerate}
        \item \term{FlashAttention}: naive attention writes the \shape{[B, H, T, T]} score matrix to HBM (quadratic memory + slow HBM round-trips). FlashAttention tiles \code{Q, K, V} into on-chip SRAM ($\sim$20 MB on H100), computes exact attention in one fused kernel, and never materializes \shape{[T, T]} in HBM $\rightarrow$ linear memory and 2--4$\times$ speedup.
        \item \term{Activation Checkpointing (Recomputation)}: discard intra-layer activations during forward and keep only layer boundaries; recompute one layer's forward pass right before its backward pass. Trades \term{$\sim$33\% extra compute} for \term{5--10$\times$ less activation HBM}.
    \end{enumerate}
\end{itemize}

% -----------------------------------------------------------------------------
\section{Precision Engineering: BF16 vs.\ DeepSeek-V3's FP8}

\begin{itemize}
    \item \term{NVIDIA H100 dense Tensor Core throughput}:
    \begin{itemize}
        \item \code{FP32}: $\sim$67 TFLOPS.
        \item \code{BF16}: $\sim$989 TFLOPS ($\sim$15$\times$ faster, half the bytes).
        \item \code{FP8}: $\sim$1,979 TFLOPS ($\sim$30$\times$ faster, quarter of the bytes).
    \end{itemize}
    \item \term{Standard recipe (Qwen2.5)}:
    \begin{itemize}
        \item Forward/backward GEMMs in \code{BF16} (matches FP32 exponent range $\rightarrow$ avoids overflow and needs no loss scaling, unlike FP16).
        \item Master weights, optimizer states, and gradient accumulation in \code{FP32}.
    \end{itemize}
    \item \term{DeepSeek-V3 FP8 mixed-precision recipe}:
    \begin{itemize}
        \item Runs almost all GEMMs in \code{FP8 (E4M3)}, accumulating partial sums in \code{FP32} inside Tensor Cores.
        \item \term{Why naive FP8 fails}: outlier channels in activations force a huge per-tensor scale factor, crushing normal values to zero.
        \item \term{Fine-grained block scaling}: independent scale factors per \code{1 $\times$ 128} activation tile and per \code{128 $\times$ 128} weight block.
        \item Sensitive ops kept in BF16/FP32: embeddings, LM head, MoE gating router, RMSNorm, attention.
        \item Optimizer moments (\code{m}, \code{v}) stored in \code{BF16} instead of FP32 to shrink static state further (master weights stay FP32).
    \end{itemize}
\end{itemize}

% -----------------------------------------------------------------------------
\section{Scaling Beyond One GPU: The 5 Parallelism Axes}

Think of every distributed training strategy by \term{what is split}, \term{what collective runs}, and \term{which physical wire it requires}:

\par\vspace{4pt}\noindent\begin{minipage}{\linewidth}
\small
\begin{tabularx}{\linewidth}{@{}l Y Y Y@{}}
\toprule
\term{Axis} & \term{What is split} & \term{Collective / Wire} & \term{Bottleneck \& Scope} \\
\midrule
\term{DP} & Batch \code{B} & \code{AllReduce} grads (IB / Eth) & Replicates full \code{16P} state on every GPU; doesn't fix OOM \\
\term{ZeRO / FSDP} & Batch \code{B} + sharded opt (Z1), grads (Z2), weights (Z3) & \code{AllGather} weights + \code{ReduceScatter} grads (IB) & Divides static HBM by \code{N\_dp}; Z1 cheap, Z3 adds weight traffic \\
\term{TP} & Matrix columns/rows (\code{D}, \code{D\_ff}, heads) & 2$\times$ \code{AllReduce} \term{per layer} (NVLink) & Strictly \term{intra-node} (8 GPUs); blocks compute on slow links \\
\term{PP} & Layers \code{L} across stages & Point-to-point \code{Send/Recv} activations (IB) & Tiny traffic, cross-node friendly; suffers from \term{pipeline bubbles} \\
\term{EP} (MoE) & Different experts on different GPUs & 2$\times$ \code{All-to-All} per MoE layer (dispatch + combine) & Cross-node bandwidth + \term{stragglers} under unbalanced routing \\
\bottomrule
\end{tabularx}
\end{minipage}\par\vspace{4pt}

\begin{itemize}
    \item \term{Hardware hierarchy rule}:
    \begin{itemize}
        \item Fastest wire (\term{NVLink}, 900 GB/s on H100 SXM): place \term{TP} (dense) or local \term{EP}.
        \item Slower wire (\term{InfiniBand}, $\sim$50 GB/s per GPU): place \term{PP}, \term{DP / ZeRO-1}, and overlapped \term{EP}.
    \end{itemize}
    \item \term{Context Parallelism (CP)} (for \code{T = 64K--128K}): splits sequence length \code{T} across GPUs and circulates \code{K}/\code{V} blocks in a ring (Ring Attention) so no GPU holds the full sequence activations.
\end{itemize}

% -----------------------------------------------------------------------------
\section{Running Examples Side-by-Side: Qwen2.5-72B vs.\ DeepSeek-V3}

\subsection{Qwen2.5-72B (Classic Dense 3D Parallelism)}
\begin{itemize}
    \item \term{Layout}: \code{TP} (intra-node) + \code{PP} (inter-node) + \code{ZeRO-1 DP}.
    \item \term{Execution profile}: every token executes identical dense GEMMs $\rightarrow$ balanced load, high Tensor Core utilization, predictable collectives.
    \item \term{Compute cost}: scales with total parameters ($\sim$6 $\times$ P FLOPs per training token $\approx$ 436 GFLOPs/token).
\end{itemize}

\subsection{DeepSeek-V3 (671B MoE on 2,048 H800 GPUs; 14.8T tokens, \texorpdfstring{$\sim$}{\textasciitilde}2.79M GPU-hours)}
\begin{itemize}
    \item \term{Layout}: \term{16-way PP + 64-way EP (across 8 nodes) + ZeRO-1 DP}, and \term{TP = 1 (no Tensor Parallelism)}.
    \item \term{Why drop TP completely?} EP already drives heavy cross-node \code{All-to-All} traffic; removing TP eliminates blocking per-layer \code{AllReduce} calls and simplifies local GEMMs.
    \item \term{DualPipe (overlapping compute and communication)}:
    \begin{itemize}
        \item Feeds micro-batches from \term{both ends} of the pipeline simultaneously.
        \item Each pipeline step overlaps one micro-batch's compute (Attention + MoE GEMMs) with another micro-batch's communication (\code{All-to-All} dispatch/combine + PP \code{Send/Recv}).
        \item Hides nearly all cross-node communication latency even on bandwidth-capped H800 links.
    \end{itemize}
    \item \term{Auxiliary-loss-free MoE load balancing}:
    \begin{itemize}
        \item Routing collapse overloads one GPU (straggler) while others sit idle.
        \item Instead of an auxiliary loss penalty (which degrades perplexity), DeepSeek adds a small dynamic \term{bias} to each expert's routing score, nudged up/down after each step based on load.
    \end{itemize}
    \item \term{Multi-Token Prediction (MTP)}:
    \begin{itemize}
        \item Trains 1 extra lightweight head to predict \code{t+2} alongside \code{t+1}.
        \item Denser training signal at tiny compute overhead, and reusable as a \term{free speculative decoding drafter} at inference.
    \end{itemize}
\end{itemize}

% -----------------------------------------------------------------------------
\section{Training Decisions That Lock In Inference (Bridge to Chapter 2)}

Training outputs a static checkpoint (weights only --- gradients and AdamW states are discarded). Every architectural decision made above permanently governs serving cost:

\begin{itemize}
    \item \term{Attention design $\rightarrow$ KV cache capacity wall}:
    \begin{itemize}
        \item \term{Qwen2.5-72B} (GQA, 8 KV heads, 80 layers): \code{2 $\times$ 8 $\times$ 128 $\times$ 80 $\times$ 2B} = \term{320 KB / token} (BF16).
        \item \term{DeepSeek-V3} (MLA, latent \code{512 + 64}, 61 layers): \code{576 $\times$ 61 $\times$ 2B} = \term{$\sim$70 KB / token} (BF16).
        \item Despite \term{9.2$\times$ more total parameters}, DeepSeek-V3 needs \term{4.5$\times$ less KV cache per token}.
    \end{itemize}
    \item \term{Dense vs.\ MoE $\rightarrow$ memory footprint vs.\ compute}:
    \begin{itemize}
        \item Dense 72B: $\sim$145 GB weights (BF16) $\rightarrow$ fits on 2$\times$ H100; spends 72B-worth of compute per token.
        \item MoE 671B: $\sim$671 GB weights (FP8) $\rightarrow$ needs an entire 8$\times$ H100 node just to hold weights in HBM, even though only 37B active params compute per token.
    \end{itemize}
    \item \term{Post-training (SFT \& RL / GRPO) --- where training becomes inference}:
    \begin{itemize}
        \item \term{SFT}: identical training loop on curated instruction/CoT pairs; loss masked on prompt tokens.
        \item \term{GRPO (DeepSeek-R1)}: sample \code{G} completions per prompt (e.g., \code{G = 16}), score with deterministic verifiers (math/code), normalize rewards within the group $\rightarrow$ \term{no critic/value model} ($\sim$50\% less memory than PPO).
        \item \term{Systems punchline}: in RL post-training, \term{70--85\% of wall-clock GPU time is spent generating rollouts} --- i.e., running an autoregressive inference engine at scale.
    \end{itemize}
\end{itemize}
